\section{Perspectives}
\label{sec:future-work}

\subsection{Toward verified proof-obligation generation}

A semantics of B machines and refinement, together with a verified
proof-obligation generator, would connect the term-level models developed in
this thesis to the correctness of complete B developments. Implementing this
generator in Lean would directly benefit \barrel{}: proof obligations could be
produced over its shallow embedding, while Lean automation could discharge
routine well-definedness conditions as they are generated. Establishing the
semantic correspondence of the shallow embedding would complement the kernel's
guarantee about completed proofs and further reduce the trusted base. The same
infrastructure could serve the automated route and support an embedded B
language for constructing and verifying machines directly in Lean.

\subsection{Strengthening interactive proof}

Broader language coverage in \barrel{} must be accompanied by a richer lemma
library. Proof-producing procedures for recurring set, relational, invariant,
and refinement arguments could increase automation without enlarging the
trusted base. Together with proof-obligation generation tailored to Lean, this
would let the system exploit well-definedness information earlier and present
users with fewer routine side goals.

\barrel{} already combines automatic discharge of routine obligations with
interactive proof of the remaining ones. This architecture is also a natural
foundation for AI-assisted reasoning: an AI assistant may propose invariants,
lemmas, or proof scripts, while Lean's kernel remains responsible for
acceptance.

These directions strengthen the same principle developed throughout this
thesis: powerful automation can be incorporated into a reasoning workflow while
the scope and basis of each accepted result remain explicit.
